% Standalone problem document, generated from the adjacent README.md.
% Compile with XeLaTeX. Mathematical content is maintained in Markdown.
\documentclass[11pt,a4paper]{article}
\usepackage[fontsize=10.5pt]{scrextend}
\usepackage[margin=22mm,headheight=15pt,headsep=9mm,footskip=11mm]{geometry}
\usepackage{fontspec}
\setmainfont{texgyrepagella}[Extension=.otf,UprightFont=*-regular,BoldFont=*-bold,ItalicFont=*-italic,BoldItalicFont=*-bolditalic]
\setsansfont{texgyreheros}[Extension=.otf,UprightFont=*-regular,BoldFont=*-bold,ItalicFont=*-italic,BoldItalicFont=*-bolditalic]
\setmonofont{texgyrecursor}[Extension=.otf,UprightFont=*-regular,BoldFont=*-bold,ItalicFont=*-italic,BoldItalicFont=*-bolditalic,Scale=MatchLowercase]
\usepackage{amsmath,amssymb}
\usepackage{unicode-math}
\setmathfont{texgyrepagella-math.otf}
\usepackage{xcolor,microtype,enumitem,needspace,fancyhdr}
\usepackage{bookmark}
\definecolor{ink}{HTML}{17344B}
\definecolor{accent}{HTML}{197D80}
\definecolor{muted}{HTML}{596773}
\hypersetup{colorlinks=true,linkcolor=accent,urlcolor=accent,pdftitle={IE-11 - The
exact fifth complete-pivoting growth
factor},pdfauthor={Open Problems in Numerical Linear Algebra}}
\urlstyle{same}
\setlength{\parindent}{0pt}
\setlength{\parskip}{5pt plus 1pt minus 1pt}
\linespread{1.04}
\setlength{\emergencystretch}{3em}
\widowpenalty=10000
\clubpenalty=10000
\displaywidowpenalty=10000
\allowdisplaybreaks[1]
\setlist{nosep,leftmargin=1.5em}
\providecommand{\tightlist}{\setlength{\itemsep}{0pt}\setlength{\parskip}{0pt}}
\providecommand{\pandocbounded}[1]{#1}
\pagestyle{fancy}
\fancyhf{}
\fancyhead[L]{\sffamily\footnotesize\color{muted}OPEN PROBLEMS IN NUMERICAL LINEAR ALGEBRA}
\fancyhead[R]{\sffamily\bfseries\footnotesize\color{accent}IE-11}
\fancyfoot[L]{\parbox[b]{0.9\textwidth}{\sffamily\fontsize{8}{10}\selectfont\color{muted}Editorial ratings; a bounded search cannot certify that a problem remains unsolved.\\Literature check: 2026-09-08}}
\fancyfoot[R]{\sffamily\footnotesize\color{muted}\thepage}
\renewcommand{\headrulewidth}{0.3pt}
\renewcommand{\headrule}{\hbox to\headwidth{\color{accent}\leaders\hrule height \headrulewidth\hfill}}
\makeatletter
\renewcommand{\section}{\Needspace{5\baselineskip}\@startsection{section}{1}{0pt}{12pt plus 2pt minus 2pt}{4pt}{\sffamily\bfseries\large\color{ink}}}
\renewcommand{\subsection}{\Needspace{4\baselineskip}\@startsection{subsection}{2}{0pt}{9pt plus 1pt minus 1pt}{3pt}{\sffamily\bfseries\normalsize\color{ink}}}
\makeatother
\setcounter{secnumdepth}{0}
\begin{document}
{\sffamily\small\color{muted}Linear systems and elimination\par}
\vspace{3pt}
{\raggedright\hyphenpenalty=10000\exhyphenpenalty=10000\sffamily\bfseries\fontsize{21}{24}\selectfont\color{ink}The
exact fifth complete-pivoting growth factor\par}
\vspace{7pt}
{\sffamily\small\textbf{Difficulty:} challenging\quad\textbf{Importance:} interesting
to specialist\par}
\vspace{4pt}
{\color{accent}\hrule height 0.8pt}
\vspace{6pt}
\textbf{Status:} open.

\subsection{Context and notation}\label{context-and-notation}

All elimination in this problem is in exact arithmetic. Write
\(\|A\|_{\max}=\max_{ij}|a_{ij}|\). A pivoting path creates successive
active Schur complements \(S_1=A,S_2,\ldots,S_n\), with row/column
permutations as appropriate. Its element-growth factor is

\[
\rho(A)=\frac{\max_{1\leq j\leq n}\|S_j\|_{\max}}{\|A\|_{\max}}.
\]

Partial pivoting chooses a largest-magnitude entry in the active first
column; complete pivoting chooses one anywhere in the active matrix. If
ties occur, a universal statement includes every admissible tie choice;
a supremum includes all admissible paths. These conventions remove
implementation-dependent ambiguity. Matrices are nonsingular unless
otherwise stated.

\subsection{Problem statement}\label{problem-statement}

Let

\[
g_5=\sup\{\rho_{\mathrm{CP}}(A):A\in\mathbb R^{5\times5}\text{ nonsingular}\}.
\]

Let \(\alpha\) be the unique real root in \((4,5)\) of the integer
polynomial \(P_5\) displayed in Equation (2.15) of the primary reference
below. This definition specifies an exact algebraic number; numerically
\(\alpha\approx4.1325170786324728542\). Prove or disprove

\[
g_5=\alpha.
\]

All complete-pivoting paths are included in the supremum. The lower
bound \(g_5\geq\alpha\) is already established in the cited paper. The
unresolved part is global optimality over all allowable patterns of
active constraints, not exact evaluation of a particular numerical
example. Order five is independently singled out by the literature as
the first unresolved dimension; this catalog does not split the same
question into entries for every subsequent order.

\subsection{Reference}\label{reference}

Chen, Edelman, and Urschel,
\href{https://arxiv.org/html/2602.20390v1}{\emph{The largest 5th pivot
may be the root of a 61st degree polynomial}}, February 2026, Equation
(2.15), Conjecture 2.1, Theorem 2.2, and Theorem 3.5. The current
rigorous interval is \(\alpha\leq g_5\leq4.84\).

\subsection{Status check}\label{status-check}

Searches for the paper title and
\textit{complete pivoting maximum five 2026} found no global proof.
Shah--Urschel's August asymptotic result leaves this exact
finite-dimensional optimum undetermined.
\end{document}
